\begin{table}[htbp]
\centering
\caption{Second-Entrant Hazard: What It Responds To}
\label{tab:a8_mechanism}
\begin{tabular}{lcc}
\toprule
Specification & Hazard ratio & $p$-value \\
\midrule
Distance to nonlegal border (main measure) & $0.70$ & $0.003$ \\
Captive nonlegal population within 70mi & $1.02$ & $0.111$ \\
\midrule
Placebo: distance to nearest county in \emph{any} state & $0.92$ & $0.478$ \\
\quad Nonlegal-border distance, raced against the placebo & $0.64$ & $0.003$ \\
\quad Placebo distance, raced against nonlegal-border distance & $1.15$ & $0.308$ \\
\bottomrule
\end{tabular}
\begin{minipage}{0.95\textwidth}
\vspace{0.3em}
\footnotesize
\textit{Notes:} All rows are the rank-2 (1st$\to$2nd entrant) specification from
Table~\ref{tab:a8_hazard_rank}. The effect responds to \emph{proximity} to a nonlegal jurisdiction, survives
a placebo test against generic (legal-status-blind) border proximity, but does not scale with the
\emph{size} of the population behind that border -- consistent with a travel-cost threshold rather than a
continuously demand-elastic response, formalized in Section~\ref{sec:model}.
\end{minipage}
\end{table}
